\documentclass[12pt,a4paper]{article}

% Todo notes
\usepackage{etoolbox}
\newtoggle{todos}
\toggletrue{todos}
\togglefalse{todos}

\usepackage[top=2.5cm, bottom=2.5cm, left=2.5cm, right=2.5cm]{geometry} % Sets the size of the document margins
\usepackage[utf8]{inputenc}  % For handling input of special character
\usepackage[T1]{fontenc}  % For better encoding of special characters
% \usepackage[swedish]{babel}  % Uncomment for documents in Swedish
\usepackage{microtype}  % For improved typesetting
\usepackage[hyphens]{url}  % Improved hyphenation of URLs
\usepackage{hyperref}  % For clickable cross-references in the text
\DeclareUrlCommand\email{}  % Formatting of (clickable) email-addresses
\usepackage{enumitem}  % More formatting options for lists
\setlist{topsep=1ex,itemsep=0.5ex,parsep=0pt,partopsep=0pt}  % this sets the vertical spacing between list items and surrounding paragraphs
\usepackage[backend=bibtex,style=ieee,natbib=true]{biblatex}
\addbibresource{references.bib}
\usepackage{multirow}  % Required for table cells that span several rows
\usepackage{longtable}  % Necessary for tables that should span several pages
\usepackage{booktabs}  % Addtl. formatting options for tables
%
\usepackage{listings}  % For pretty printing of source code
%
\iftoggle{todos}{%
  % include the blue comment boxes
  \usepackage[color=blue!10,textsize=footnotesize,textwidth=25mm]{todonotes}
}{%
  % exclude the blue comment boxes
  \usepackage[disable]{todonotes}
}

\usepackage{float}
\usepackage{tabularx}



\title{Git For Dummies - Version 1.0.1\\\vspace{1mm}}
\author{\textsc{Linus Holmström}}
\date{\today}

\begin{document}


\maketitle

\vspace*{-12mm}

\begin{small}
\begin{longtable}{|p{13.3cm}|}
\hline
\textbf{Topics not yet covered / Planned changes} \\
\hline
\begin{itemize}
    % Not included in V1
    \item Git CLI
    \item git-flow-next
    \item Tagging commits (releases)
    \item Submodules

    % TODO
    \item Merge conflicts
\end{itemize} \\ \hline
\end{longtable}
\end{small}

\setcounter{table}{0}  % The tables on the front page shouldn't be counted


\setcounter{tocdepth}{3}
\tableofcontents

\newpage

\section{Preface}
    In this document, some features of Git and Gitflow, as well as ways to work with them, are explained. While some aspects are true in general scenarios, some aspects may be explained as true despite them only being so for this specific project. Before applying certain ways of working outside of this project, ensure that they are correct for your specific use case.

\section{Git - "Goddamn idiotic truckload of shit"}
\label{sec:git}
    Simply put, Git is a version control system. By using "commits", it tracks changes to the project ("repository"), and a sequence of commits describes how a project's files have changed over time. A git repository contains one or more branches. A branch represents a line of development, allowing developers to work on separate changes without interfering with each other. The default branch is commonly called \verb|main|.
    
    The most basic way git works is via branches, as mentioned above, and commits. Commits are collections of changes to the files and file structure of the repository. Commits can only contain changes to files, and will not include empty directories. To create a commit, you must first \textbf{stage} changes, which are then grouped with a commit message. You can create several commits on your local system before \textbf{pushing} changes in your local repository to the remote (the repository stored on GitHub). How staging works varies depending on the software used to interact with Git. To ensure clarity, commit messages must communicate clearly what changed. Read more about commits in section \ref{sec:commits}.
    
    One important distinction to make is that Git is a system. In turn, we are using GitHub, which is a repository host. Later, in section \ref{sec:pr}, we will explore pull requests. Pull requests are a feature not inherently present in Git, and exist because of our choice of repository host. There are several more such features not covered in this document.

\section{Gitflow 101}
    \label{sec:gitflow}
    As explained in section \ref{sec:git}, git is based on branches. Gitflow is a workflow used to manage branches in a project and is based on several branch types, most importantly:
    \begin{itemize}
        \item \verb|main|
        \item \verb|dev|
        \item Feature branches
        \item Bugfix branches
        \item Release branches
    \end{itemize}
    
    Any changes should always be done on at least a feature branch, \textbf{not} on \verb|main| or \verb|dev|. Feature branches and other types are explained below.

    \subsection{Branches in Gitflow}
    \label{subsec:gitflow_branches}
        
        \subsubsection{The main branch}
            The main branch, not surprisingly named \verb|main|, is \textbf{the} branch, at least in our case. In the case of Project Thalassa and Droplet, \verb|main| acts as the release branch, always containing a buildable, runnable version of the project. Making changes to \verb|main| should be done sparingly and cautiously, ensuring that a change to the branch is thoroughly tested to avoid major issues. Changes to \verb|main| always originate from a release branch, except for the rare case of hotfixes.
        
        \subsubsection{The development branch}
            The development branch, \verb|dev|, contains the most recent version of the project with finished features. The development branch must always build and run, but, unlike main, it is acceptable for bugs and unfinished features to be present on \verb|dev|. Changes to \verb|dev| always originate from feature, bugfix, or release branches.
        
        \subsubsection{Feature branches}
            Feature branches are temporary branches, with a lifetime limited to the time needed to implement the feature they are intended for. Feature branches are always used for new functionality in the project. Once a feature is merged into the \verb|dev| branch, its corresponding feature branch should be deleted. Names of feature branches are always \verb|feature/<name-of-feature>|, and are always based on the latest \verb|dev| commit.
        
        \subsubsection{Bugfix branches}
            Bugfix branches, like feature branches, are temporary branches. The purpose of a bugfix branch is to fix a specific bug within the project. Bugfix branches are always \verb|bugfix/<name-of-bugfix>|.
        
        \subsubsection{Release branches}
            Release branches are created to prepare a change to \verb|main|. Like feature and bugfix branches, release branches are deleted once their purpose is fulfilled. The purpose of release branches is to apply any changes and final tweaks, and to test the new release thoroughly before performing the release. Once ready, release branches should be merged into \textbf{both} \verb|main| and \verb|dev|, the latter to establish any fixes and tweaks made during the release candidate work. Release branches are always named \verb|release/<version-number>|.
        
        \subsubsection{Hotfix branches}
            Hotfix branches should be extremely rare. Their sole purpose is to fix urgent, project-breaking bugs. Like most branches, they live only for the duration needed to implement the changes they are meant for. Hotfix branches always originate from \verb|main|, and should always be merged into \textbf{both} \verb|main| and \verb|dev|. Hotfix branches are always named \verb|hotfix/<name-of-hotfix>|.

    \subsection{Workflow in Project Thalassa and Droplet}
    \label{subsec:workflow}
        For Project Thalassa and Droplet, the workflow for a branch (branch A) that is ready to start the process of being merged into another branch (branch B) is as follows:
        \begin{enumerate}
            \item Ensure that branch A at least has the latest changes from the remote copy.
            \item Ensure that the local copy of branch B is up to date, \textbf{or} ensure that you will be referencing the remote copy.
            \item Using your client of choice, merge branch B into branch A. This allows you to handle merge conflicts locally.
            \item Handle any merge conflicts and push your changes to the remote.
            \item Verify that branch A builds and runs correctly after the merge.
            \item Create a pull request to merge the now up-to-date branch A into branch B.
            \item If branch A has served its purpose, delete it once merged.
        \end{enumerate}
    
\section{Pull requests and you}
\label{sec:pr}
    In the project, making changes to the \verb|dev| and \verb|main| branches requires pull requests. A pull request (PR) is a request to merge changes from one branch into another. Pull requests on \verb|dev| require approval by one other person, and changes on \verb|main| require approval by three others. Making changes to any other branches is not regulated and can be done freely.

    \subsection{Creating a pull request}
        If you have changes in your branch, let's call it \verb|feature/jumping|, that you are satisfied with and want added into the \verb|dev| branch, you need to create a Pull Request. A PR is easiest to create via the GitHub webpage.
    
        \begin{enumerate}
            \item Go to the repository where your branch is located.
            \item Click on "Pull requests".
            \item Click on "New pull request".
            \item Select \verb|dev| as "base" and your branch as "compare".
            \item Write your title and description.
            \item Click "Create pull request".
        \end{enumerate}
        
        Once the PR is created, you can either ask others to check it out and approve it, or you can assign others to it directly.
    
        Before creating a pull request, you should merge the destination branch into your branch. See section \ref{subsec:branch_merge} for details on merging, and section \ref{subsec:workflow} for details on the process of merging.
        
    \subsection{Reviewing a pull request}
        Pull requests must be approved by other collaborators before it is possible to merge the changes.
    
        \begin{enumerate}
            \item Go to the repository where the pull request is located.
            \item Click on "Pull requests".
            \item Click on the pull request to be reviewed.
            \item Read the title and description for any necessary context.
            \item Click on "Files changed".
            \item Study the changes made. For ease, you may mark files as viewed on the right-hand side.
            \item If you see changes that should be made, you may add comments on the entire file on the right-hand side, or on a specific line by clicking the plus icon when hovering over it.
            \item When satisfied, click the green button in the top right and select the option that applies. Note that "Approve" means that no changes are expected or suggested, and should only be used if things look good to merge.
        \end{enumerate}

        
    
        When a PR is fully approved, anyone is able to merge the changes into the base branch, finalising the merge. All branches that are not \verb|main| and \verb|dev| should be deleted once merged; see section \ref{subsec:gitflow_branches}.
    
        \subsubsection{Review checklist}
            When reviewing a PR, consider the following:
            \begin{itemize}
                \item Do the changes do what the pull request claims?
                \item Is there anything that is obviously incorrect?
                \item Is the code understandable?
                \item Does the code follow the agreed-upon standards?
                \item Do the changes introduce unnecessary complexity?
                \item Are there debugging leftovers that serve no purpose?
            \end{itemize}
    
        \subsubsection{Testing a pull request}
            Pull requests must always be \textbf{tested} unless obviously unnecessary, e.g., if documentation has changed. This can be done by pressing "Code" in the top right of the pull request page and using the option that applies ("Checkout with GitHub Desktop" or "Checkout with GitHub CLI"). Once checked out, ensure functionality and that the project compiles and runs. When doing this, you will be "checking out" a version of the project, which means that your changes must be saved and that you then must return to your branch. See section \ref{subsec:branch_switch} for more details.


\section{Commits}
\label{sec:commits}
    Commits represent snapshots of the repository, and each commit keeps a reference to its predecessor(s). Using a sequence of commits, it is possible for Git to construct the repo. A commit is always accompanied by a commit message and is always tagged with a unique identifier.

    \subsection{Writing good commit messages}
    \label{subsec:commit_messages}
        Writing good commit messages is somewhat of an art form. Their most important purposes are to speed up the reviewing process, help others write good release notes, and help future maintainers understand why certain changes were made \cite{hogberg_writing_2023}. Conventional Commits 1.0.0 \cite{noauthor_conventional_nodate} provides a good specification for how commit messages should be formatted, and commit messages in Thalassa and Droplet should follow Conventional Commits 1.0.0. While compliance is not automatically enforced through any systems, contributors are expected to follow the specification.

        Conventional Commits 1.0.0 states that commit messages should be structured as follows, with a short explanation of the elements below. Readers are encouraged to study the specification on \textcite{noauthor_conventional_nodate}.

        \begin{verbatim}
            <type>[optional scope]: <description>

            [optional body]

            [optional footer(s)]
        \end{verbatim}

        % Please add the following required packages to your document preamble:
        \begin{table}[H]
        \centering
        \begin{tabularx}{\linewidth}{l|>{\raggedright\arraybackslash}X}
        \textbf{Element} & \textbf{Explanation}                                                                                                                         \\ \hline
        type             & The kind of change, e.g. \verb|feat| or \verb|fix|.                                                                                          \\ \hline
        scope            & Optional indication of which system is affected, e.g. renderer, asset manager.  If a scope is included, it should be encased in parentheses. \\ \hline
        description      & A summary of the change.                                                                                                               \\ \hline
        body             & Optional additional context, describing what changed and why.                                                                                \\ \hline
        footer           & Optional metadata, such as references to issues or breaking changes.                                                                        
        \end{tabularx}
        \end{table}

        \subsubsection{Example commit messages}
            Below, some example commit messages are shown to contextualize section \ref{subsec:commit_messages}.
            
            \verb|feat(renderer): add support for textured meshes| \vspace{2mm}
            
            \verb|fix(physics): correct collision response on slopes| \vspace{2mm}
            
            \verb|docs: update Git workflow documentation| \vspace{2mm}
            
            \begin{verbatim}
    fix(renderer): correct image layout transition
    
    Transition sampled images to the appropriate layout
    before they are accessed by shaders.
    
    This fixes validation errors occurring during rendering.
            \end{verbatim}

    \subsection{Sizing a commit}
    \label{subsec:commit_sizes}
        It is easy to get carried away, dumping all changes in one large commit. Doing this is not recommended, however. With large commits, it becomes hard to review changes, pinpoint sources of bugs, and revert breaking changes. Instead, commits should represent logically self-contained changes. A commit should not be unnecessarily small, either. Many small commits that together implement or modify one feature may be better suited in one commit, even if they affect multiple files.
        
        A good rule of thumb is that if it is hard to summarize a commit, you want to rethink the size and split the changes into several commits.

\section{Branch management}
\label{sec:branch_management}
    Branch management can become tricky if you don't understand the subject at all. This section will cover the most basic branch concepts needed to work with them in the project.

    \subsection{Switching branches}
    \label{subsec:branch_switch}
        You can always "check out" a branch. Checking out a branch simply means switching which branch is present on your machine. To make the switching process easy, it is best not to have any uncommitted changes present (see section \ref{sec:commits}), though it is possible to "stash" changes. You will be prompted to do this, so there is no worry about losing work.

    \subsection{Rebasing branches}
    \label{subsec:branch_rebase}
        Since Git works with a timeline, this timeline might get jumbled if changes have been made on a branch you intend to merge yours into. Say you created a feature branch from \verb|dev| when commit 15 was the most recent one. When you are finished with your feature, \verb|dev| might now be on commit 20, and the branches have "diverged". In some cases, Git can automatically merge changes, but sometimes it can't decide which changes to keep or discard, and you get a merge conflict. To avoid issues, it is recommended to "rebase" your branch. Rebasing a branch means taking all changes made in the branch and attempting to apply them with a different commit as the base. Most of the time, a rebase will create a linear timeline. How you perform rebases varies depending on your client of choice.

        In Thalassa and Droplet, developers are \textbf{not} expected to rebase their changes before creating pull requests.

    \subsection{Merging branches}
    \label{subsec:branch_merge}
        Merging branches is another way of dealing with the scenario described in section \ref{subsec:branch_rebase}. The key difference between rebasing and merging is that, with merging, you take the \verb|head|s of two branches and combine the two snapshots of the repository, effectively merging two timelines into one. When merging branch A into branch B, git combines their histories. If the branches have diverged, git normally creates a merge commit, which contains the combined result. Otherwise, git performs a fast-forward merge, which does not create a merge commit. Unlike rebasing, merging preserves the existing commit history, and in the case of a merge commit, the history becomes non-linear and reflects the actual development history. This is not always preferable and depends on the project. How you perform merges varies depending on your client.
        

\small
\printbibliography

\end{document}
